\FigureLayoutDeclare{img/1963/national/q3_fig1.png}{6.0cm}{0bp 0bp 0bp 0bp}
\FigureLayoutDeclare{img/1963/national/q3_solution_fig1.png}{4.0cm}{11bp 7bp 13bp 11bp}
\FigureLayoutDeclare{img/1963/national/q4_solution_fig1.png}{4.0cm}{0bp 0bp 0bp 0bp}
\FigureLayoutDeclare{img/1963/national/q9_fig1.png}{6.0cm}{0bp 0bp 0bp 0bp}
\FigureLayoutDeclare{img/1963/national/q9_solution_fig1.png}{4.2cm}{0bp 0bp 0bp 0bp}

\chapter{1963年全国卷}

\section{解答题}

\begin{problem}
已知 \(\tan\theta=\sqrt2\)，求 \(\dfrac{\cos\theta+\sin\theta}{\cos\theta-\sin\theta}\) 的值.
\end{problem}

\begin{solution}
因为 \(\tan\theta=\sqrt2\)，所以 \(\cos\theta\neq0\)，且 \(1-\tan\theta\neq0\). 以 \(\cos\theta\) 除分式的分子和分母，得
\[
\frac{\cos\theta+\sin\theta}{\cos\theta-\sin\theta}
=\frac{1+\tan\theta}{1-\tan\theta}
=\frac{1+\sqrt2}{1-\sqrt2}
=\frac{\left( 1+\sqrt2 \right)^2}{\left( 1-\sqrt2 \right)\left( 1+\sqrt2 \right)}
\]
\[
\frac{1+2\sqrt2+2}{-1}=-3-2\sqrt2
\]

因此原式的值为 \(-3-2\sqrt2\).
\end{solution}

\begin{problem}
\begin{enumerate}
\item 求出复数 \(1+\sqrt3\i\) 的模数和辐角.
\item 在直角坐标系 \(XOY\) 所在的平面内，以 \(A\) 点表示复数 \(1+\sqrt3\i\)，把 \(OA\) 绕着 \(O\) 点按反时针方向旋转 \(150^\circ\)，设 \(A\) 点到达的位置为 \(B\)，写出 \(B\) 点所表示的复数的代数式.
\end{enumerate}
\end{problem}

\begin{solution}
\begin{enumerate}
\item 设 \(z=1+\sqrt3\i\). 由复数模的定义，
\(|z|=\sqrt{1^2+\left(\sqrt3\right)^2}=2\).
因为实部为正、虚部为正，且若主值为 \(\alpha\)，则 \(\tan\alpha=\sqrt3\)，所以 \(z\) 位于第一象限，辐角的主值为 \(60^\circ\). 因而 \(z\) 的全部辐角为 \(60^\circ+360^\circ k\)，其中 \(k\in\Z\).

\item 旋转不改变模，且逆时针旋转使辐角增加 \(150^\circ\)，所以点 \(B\) 表示的复数的模为 \(2\)，辐角主值为
\(60^\circ+150^\circ=210^\circ\).
因此点 \(B\) 表示的复数为
\[
2\left(\cos210^\circ+\i\sin210^\circ\right)
=2\left(-\frac{\sqrt3}{2}-\frac12\i\right)
=-\sqrt3-\i
\]

\end{enumerate}
\solutionfigure{\bitmapfigure[width=3cm]{img/1963/national/q2_solution_fig1.png}}
\end{solution}

\begin{problem}
如图，\(AB\) 为半圆的直径，\(CD\perp AB\). 已知 \(AB=1\)，\(AC:CB=4:1\)，求 \(CD\).

\bitmapfigure[width=6.0cm]{img/1963/national/q3_fig1.png}
\end{problem}

\begin{solution}
因为 \(AB=1\)，且 \(AC:CB=4:1\)，所以 \(AC=\frac45\)，\(CB=\frac15\).
由于 \(AB\) 是半圆的直径，圆周角 \(\angle ADB=90^\circ\). 又 \(CD\perp AB\)，所以 \(CD\) 是直角三角形 \(ADB\) 斜边 \(AB\) 上的高.

在直角三角形 \(ACD\) 和 \(DCB\) 中，\(\angle ACD=\angle DCB=90^\circ\). 因为 \(C\) 在线段 \(AB\) 上且 \(DC\) 位于直角 \(\angle ADB\) 内，所以 \(\angle ADC+\angle CDB=90^\circ\). 又由 \(\angle ACD=90^\circ\)，有 \(\angle CAD+\angle ADC=90^\circ\)，从而 \(\angle CAD=\angle CDB\). 因此
\[
\triangle ACD\sim\triangle DCB
\]
由相似三角形对应边成比例，得
\(\frac{AC}{CD}=\frac{CD}{CB}\)，\(CD^2=AC\cdot CB =\frac45\cdot\frac15=\frac4{25}\).
因为 \(CD>0\)，所以 \(CD=\frac25\).

\solutionfigure{\bitmapfigure[width=4.0cm]{img/1963/national/q3_solution_fig1.png}}
\end{solution}

\begin{problem}
从二面角内任意一点向二面角的两个面作垂线，求证这两条垂线所决定的平面垂直于二面角的棱（要求画图）.
\end{problem}

\begin{solution}
设二面角的两个面为 \(M\)、\(N\)，棱为 \(AB\)，点 \(P\) 在二面角内，且 \(PC\perp M\)、\(PD\perp N\)，其中 \(C\in M\)、\(D\in N\). 直线 \(PC\)、\(PD\) 相交于 \(P\)，所以它们确定平面 \(PCD\). 因为平面 \(PCD\) 经过直线 \(PC\)，而 \(PC\perp M\)，所以平面 \(PCD\perp M\). 同理，平面 \(PCD\perp N\).

平面 \(M\)、\(N\) 相交于棱 \(AB\). 两个相交平面同时垂直于第三个平面时，它们的交线垂直于第三个平面，因此 \(AB\perp\) 平面 \(PCD\)，即平面 \(PCD\perp AB\).

\solutionfigure{\bitmapfigure[width=4cm]{img/1963/national/q4_solution_fig1.png}}
\end{solution}

\begin{problem}
利用下列常用对数表，计算 \(23.28^{-1.1}\).

\begin{center}
\scriptsize
\setlength{\tabcolsep}{3pt}
\renewcommand{\arraystretch}{1.1}
\begin{tabular}{|c|cccccccccc|ccccccccc|}
\hline
\multicolumn{20}{|c|}{对数表} \\
\hline
\(N\) & 0 & 1 & 2 & 3 & 4 & 5 & 6 & 7 & 8 & 9 & 1 & 2 & 3 & 4 & 5 & 6 & 7 & 8 & 9 \\
\hline
23 & 3617 & 3636 & 3655 & 3674 & 3692 & 3711 & 3729 & 3747 & 3766 & 3784 & 2 & 4 & 6 & 7 & 9 & 11 & 13 & 15 & 17 \\
24 & 3802 & 3820 & 3838 & 3856 & 3874 & 3892 & 3909 & 3927 & 3945 & 3962 & 2 & 4 & 5 & 7 & 9 & 11 & 12 & 14 & 16 \\
25 & 3979 & 3997 & 4014 & 4031 & 4048 & 4065 & 4082 & 4099 & 4116 & 4133 & 2 & 3 & 5 & 7 & 9 & 10 & 12 & 14 & 15 \\
\hline
\end{tabular}

\vspace{0.8em}

\begin{tabular}{|c|cccccccccc|ccccccccc|}
\hline
\multicolumn{20}{|c|}{反对数表} \\
\hline
\(m\) & 0 & 1 & 2 & 3 & 4 & 5 & 6 & 7 & 8 & 9 & 1 & 2 & 3 & 4 & 5 & 6 & 7 & 8 & 9 \\
\hline
.48 & 3020 & 3027 & 3034 & 3041 & 3048 & 3055 & 3062 & 3069 & 3076 & 3083 & 1 & 1 & 2 & 3 & 4 & 4 & 5 & 6 & 6 \\
.49 & 3090 & 3097 & 3105 & 3112 & 3119 & 3126 & 3133 & 3141 & 3148 & 3155 & 1 & 1 & 2 & 3 & 4 & 4 & 5 & 6 & 6 \\
.50 & 3162 & 3170 & 3177 & 3184 & 3192 & 3199 & 3206 & 3214 & 3221 & 3228 & 1 & 1 & 2 & 3 & 4 & 4 & 5 & 6 & 7 \\
\hline
\end{tabular}
\end{center}
\end{problem}

\begin{solution}
由对数表，在数表的第 \(23\) 行、第 \(2\) 列查得 \(1.3655\)，再由横向的平均差表查得第 \(8\) 个平均差 \(15\)，所以
\[
\lg 23.28=1.3655+0.0015=1.3670
\]
于是
\[
\lg 23.28^{-1.1}=-1.1\lg 23.28=-1.1\times1.3670=-1.5037
=\overline{2}.4963
\]
查反对数表，\(.4963\) 在 \(.49\) 行、第 \(6\) 列的数为 \(3133\)，再加上第 \(3\) 个平均差 \(2\)，得到 \(3135\). 因而
\(10^{.4963}\approx3.135\)，\(23.28^{-1.1}=10^{\overline{2}.4963} =10^{-2}\times10^{.4963}\approx0.03135\).
\end{solution}

\begin{problem}
解方程 \(\sin 3x-\sin x+\cos 2x=0\).
\end{problem}

\begin{solution}
【法一】
\[
\sin 3x-\sin x+\cos 2x=0
\Longleftrightarrow 2\cos 2x\sin x+\cos 2x=0
\Longleftrightarrow \cos 2x\left( 2\sin x+1 \right)=0
\]

由 \(\cos 2x=0\)，得 \(2x=2n\pi\pm\frac{\pi}{2}\)，所以 \(x=n\pi\pm\frac{\pi}{4}\).

由 \(2\sin x+1=0\)，得 \(\sin x=-\frac12\)，所以 \(x=2n\pi-\frac{\pi}{6}\)，或 \(x=\left( 2n+1 \right)\pi+\frac{\pi}{6}\).

所以 \(x=n\pi\pm\frac{\pi}{4}\)，\(x=2n\pi-\frac{\pi}{6}\)，或 \(x=\left( 2n+1 \right)\pi+\frac{\pi}{6}\)（\(n\in\Z\)）.

【法二】
\[
\begin{gathered}
\sin 3x-\sin x+\cos 2x
=\left( 3\sin x-4\sin^3x \right)-\sin x+\left( 1-2\sin^2x \right)=0\\
\Longleftrightarrow 4\sin^3x+2\sin^2x-2\sin x-1
=\left( 2\sin x+1 \right)\left( 2\sin^2x-1 \right)=0
\end{gathered}
\]

由 \(2\sin x+1=0\)，得 \(\sin x=-\frac12\)，所以 \(x=2n\pi-\frac{\pi}{6}\)，或 \(x=\left( 2n+1 \right)\pi+\frac{\pi}{6}\).

由 \(2\sin^2x-1=0\)，得 \(\sin x=\pm\frac{\sqrt2}{2}\)，所以 \(x=2n\pi+\frac{\pi}{4}\)，\(x=\left( 2n+1 \right)\pi-\frac{\pi}{4}\)，\(x=2n\pi-\frac{\pi}{4}\)，或 \(x=\left( 2n+1 \right)\pi+\frac{\pi}{4}\).

所以 \(x=2n\pi-\frac{\pi}{6}\)，\(x=\left( 2n+1 \right)\pi+\frac{\pi}{6}\)，\(x=n\pi\pm\frac{\pi}{4}\)（\(n\in\Z\)）.
\end{solution}

\begin{problem}
用 \(1\)，\(2\)，\(3\)，\(4\)，\(7\)，\(9\) 组成没有重复数字的五位数，问：

\begin{enumerate}
\item 这样的五位数一共有多少个？
\item 在这些五位数中，有多少个是偶数？
\item 在这些五位数中，有多少个是 \(3\) 的倍数？
\end{enumerate}
\end{problem}

\begin{solution}
\begin{enumerate}
\item 从这六个数字中取出五个数字，共能排成 \(\mathrm A_6^5=6\times5\times4\times3\times2=720\) 个五位数.

\item 在所求的偶数中，末位必须取 \(2\)，\(4\) 这两个数字中的一个，这有两种方法. 取定末位后，再从其余五个数字中任取四个，排成其他四位，这有 \(\mathrm A_5^4=5\times4\times3\times2=120\) 种方法. 因此，共有 \(2\mathrm A_5^4=2\times120=240\) 个五位数是偶数.

\item 一个整数是不是 \(3\) 的倍数，要看它的各位数字之和是不是 \(3\) 的倍数. 这六个数字 \(1\)，\(2\)，\(3\)，\(4\)，\(7\)，\(9\) 之和是 \(26\). 取出五个数字时，若删去的数字为 \(d\)，则所取数字之和为 \(26-d\). 要使 \(26-d\) 是 \(3\) 的倍数，必须有 \(d\equiv2\pmod3\). 在这六个数字中，只有 \(2\) 满足这个条件，所以所取的五个数字必须是 \(1\)，\(3\)，\(4\)，\(7\)，\(9\). 因此，共 \(\mathrm A_5^5=5!=120\) 个五位数是 \(3\) 的倍数.
\end{enumerate}
\end{solution}

\begin{problem}
解方程组：\(\begin{cases}
x^2-2xy-y^2=1,\\
\sqrt{xy+3}=x
\end{cases}\)（限定在实数范围内）.
\end{problem}

\begin{solution}
第二式要求 \(x=\sqrt{xy+3}\geq0\). 对第二式平方，得到
\(xy+3=x^2\)，即 \(x^2-xy=3\).
将第一式乘以 \(3\)，再减去上式，得
\[
3x^2-6xy-3y^2-(x^2-xy)=3-3
\]
即
\[
2x^2-5xy-3y^2=0
\Longleftrightarrow
\left(2x+y\right)\left(x-3y\right)=0
\]

当 \(2x+y=0\) 时，\(y=-2x\). 代入 \(x^2-xy=3\)，得
\(3x^2=3\)，所以 \(x=\pm1\).
于是得到候选解 \((x,y)=(1,-2)\) 和 \((-1,2)\). 对 \((1,-2)\)，第一式左边为 \(1-2\cdot1\cdot(-2)-(-2)^2=1\)，第二式为 \(\sqrt{1}=1\)，所以 \((1,-2)\) 是解. 对 \((-1,2)\)，第一式虽也成立，但第二式为 \(\sqrt{1}=1\neq-1\)，所以舍去.

当 \(x-3y=0\) 时，\(x=3y\). 代入 \(x^2-xy=3\)，得
\(6y^2=3\)，所以 \(y=\pm\frac{\sqrt2}{2}\)，\(x=\pm\frac{3\sqrt2}{2}\).
当 \(x=\frac{3\sqrt2}{2}\)、\(y=\frac{\sqrt2}{2}\) 时，第一式左边为 \(\frac92-3-\frac12=1\)，且 \(xy+3=\frac92\)，所以第二式为 \(\sqrt{\frac92}=\frac{3\sqrt2}{2}=x\)，满足原方程组. 当 \(x=-\frac{3\sqrt2}{2}\)、\(y=-\frac{\sqrt2}{2}\) 时，第一式仍成立，但左边 \(\sqrt{xy+3}=\frac{3\sqrt2}{2}\) 为正数，而右边为负数 \(x=-\frac{3\sqrt2}{2}\)，不满足第二式.

因此，原方程组的解为
\(\left(x,y\right)=\left(1,-2\right)\)，\(\left(x,y\right)=\left(\frac{3\sqrt2}{2},\frac{\sqrt2}{2}\right)\).
\end{solution}

\begin{problem}
如图，线段 \(CD\) 与 \(\odot O\) 相交于 \(A\)，\(B\) 两点，且 \(AC=BD\)，又 \(CE\)，\(DF\) 分别与 \(\odot O\) 相切于 \(E\)，\(F\). 求证：\(\triangle OEC\cong\triangle OFD\).

\bitmapfigure[width=6.0cm]{img/1963/national/q9_fig1.png}
\end{problem}

\begin{solution}
连结 \(OA\)、\(OB\). 因为 \(OA\)、\(OB\) 都是半径，所以 \(OA=OB\)，从而在等腰三角形 \(OAB\) 中
\[
\angle OAB=\angle OBA
\]

\solutionfigure{\bitmapfigure[width=4.2cm]{img/1963/national/q9_solution_fig1.png}}

按图中点的顺序，\(C\)、\(A\)、\(B\)、\(D\) 共线，因此
\(\angle OAC=180^\circ-\angle OAB\)，\(\angle OBD=180^\circ-\angle OBA\).
由 \(\angle OAB=\angle OBA\)，可得 \(\angle OAC=\angle OBD\). 又已知 \(AC=BD\)，所以
\[
\triangle OAC\cong\triangle OBD
\]
（边角边），从而 \(OC=OD\).

因为 \(CE\)、\(DF\) 分别是 \(\odot O\) 在 \(E\)、\(F\) 处的切线，所以半径垂直切线，
\[
\angle OEC=\angle OFD=90^\circ
\]
又 \(OE=OF\) 是同一圆的半径，且已经得到 \(OC=OD\). 因此两个直角三角形的斜边及一条直角边分别相等，故
\[
\triangle OEC\cong\triangle OFD
\]
（斜边直角边）.
\end{solution}

\begin{problem}
半径为 \(1\) 的球，内切于圆锥（即直圆锥），已知圆锥的母线与底面的夹角是 \(2\theta\).

\begin{enumerate}
\item 求证圆锥的母线与底面半径的和是 \(\dfrac{2}{\tan\theta\left(1-\tan^2\theta\right)}\).
\item 求证圆锥的全面积是 \(\dfrac{2\pi}{\tan^2\theta\left(1-\tan^2\theta\right)}\).
\item 当 \(\theta\) 是什么值的时候，圆锥的全面积最小？（\(\theta\) 用反三角函数表表示）
\end{enumerate}

（图中 \(V\) 是圆锥的顶点，\(VB\) 是母线，\(O\) 是球心，\(A\) 是球和圆锥底面的切点.）

\problemresetlayout
\bitmapfigure[width=6.0cm]{img/1963/national/q10_fig1.png}
\end{problem}

\begin{solution}
\begin{enumerate}
\item 设 \(C\) 为母线 \(VB\) 与球相切的切点，连接 \(OA\)、\(OB\)、\(OC\). 因为球与底面在 \(A\) 点相切，球与母线在 \(C\) 点相切，所以
\(OA=OC=1\)，且 \(\angle OAB=\angle OCB=90^\circ\).
三角形 \(OAB\) 与 \(OCB\) 有公共斜边 \(OB\)，且有一条直角边 \(OA=OC\)，所以
\[
\triangle OAB\cong\triangle OCB
\]
从而 \(\angle OBA=\angle OBC\). 又因为 \(BC\) 在母线 \(BV\) 上，且母线与底面的夹角为 \(2\theta\)，故
\[
\angle OBA+\angle OBC=\angle ABV=2\theta
\]
因此 \(\angle OBA=\angle OBC=\theta\).

设圆锥的底面半径为 \(r\)，母线为 \(l\). 在直角三角形 \(OAB\) 中，
\(\tan\theta=\frac{OA}{AB}=\frac1r\)，\(r=\cot\theta\).
在直角三角形 \(VAB\) 中，\(AB=r\)、\(VB=l\)，且 \(\angle ABV=2\theta\)，所以
\(\cos2\theta=\frac{AB}{VB}=\frac rl\)，\(l=\frac{r}{\cos2\theta}=\frac{\cot\theta}{\cos2\theta}\).
圆锥存在时 \(0<2\theta<90^\circ\)，故 \(0<\theta<45^\circ\)，从而 \(1-\tan^2\theta>0\). 于是
\[
l+r=\cot\theta\left(\frac{1}{\cos2\theta}+1\right)
=\frac{1}{\tan\theta}\cdot\frac{1+\cos2\theta}{\cos2\theta}
=\frac{1}{\tan\theta}\cdot\frac{2\cos^2\theta}{\cos^2\theta-\sin^2\theta}
\]
\[
l+r=\frac{1}{\tan\theta}\cdot\frac{2}{1-\tan^2\theta}
=\frac{2}{\tan\theta\left(1-\tan^2\theta\right)}
\]

\item 设圆锥的全面积为 \(T\). 全面积等于侧面积与底面积之和，所以
\[
T=\pi rl+\pi r^2
=\pi r\left(l+r\right)
=\pi\cot\theta\cdot\frac{2}{\tan\theta\left(1-\tan^2\theta\right)}
\]
\[
T=\frac{2\pi}{\tan^2\theta\left(1-\tan^2\theta\right)}
\]

\item 由前一问，
\[
T=\frac{2\pi}{\tan^2\theta\left(1-\tan^2\theta\right)}
\]
在 \(0<\theta<45^\circ\) 内，令 \(u=\tan^2\theta\)，则 \(0<u<1\). 因为 \(T>0\)，要使 \(T\) 最小，只需使
\[
u(1-u)=-u^2+u=-\left(u-\frac12\right)^2+\frac14
\]
最大. 当且仅当 \(u=\frac12\) 时取得最大值，因此
\[
\tan^2\theta=\frac12
\]
又 \(0<\theta<45^\circ\)，所以 \(\tan\theta>0\)，从而
\[
\theta=\arctan\frac{\sqrt2}{2}
\]
此时圆锥的全面积最小.
\end{enumerate}
\solutionfigure{\bitmapfigure[width=3.6cm]{img/1963/national/q10_solution_fig1.png}}
\end{solution}
